% =============================================================================
% Ambient-electron exchange thruster -- arXiv preprint (v2, characterization campaign)
% Build: pdflatex main.tex  (twice for refs)
% Status: DRAFT. Measured numbers from the gated capstone frontier
% (100/200/300/350 V + slender pair, reference_results/) + minimal model mission sweeps.
% TODO markers where inputs are pending (lab experiment, author list,
% final lit refs). Characterization spokes (density, magnetized) folded
% in 2026-08-12.
% Lineage: reworked from an earlier internal draft (O14 mm / 800 V
% skeleton); predecessor results retained as supporting evidence.
% =============================================================================
\documentclass[11pt]{article}
\usepackage[margin=1in]{geometry}
\usepackage{amsmath,amssymb,graphicx,booktabs,hyperref,xcolor,rotating}
\newcommand{\todo}[1]{\textcolor{red}{[TODO: #1]}}

\title{Propellantless Drag Compensation by Ambient-Electron Exchange:\\
       a Measured Thrust--Power Frontier and a Scale-Free Feasibility Condition}
\author{R.~S.~C.\thanks{\todo{author list, affiliation, email}}}
\date{\today}

\begin{document}
\maketitle

\begin{abstract}
We describe and simulate a spacecraft thruster with no stored propellant, no
tank, and no neutralizer: a cold electron gun sealed inside a floating
conducting body emits an electron beam through a single aperture; the body
charges positive until it recollects ambient ionospheric electrons at exactly
the emitted rate. In the ideal limit the concept's economics are
parameter-free -- thrust $F = \sqrt{2m_e/e}\,I\sqrt{\mathrm{KE}}$ at beam
power $P = VI$, a geometry-free lower bound on fixed-thrust power set by
electron mass and charge alone. Taking a \O10\,mm conducting can at a
representative dayside ionospheric operating point as a measured anchor,
pre-registered and acceptance-gated particle-in-cell simulations (WarpX, RZ
electrostatic) validate that concept under real constraints, measuring the
complete thrust--power frontier across the 100--350\,V hardware
range: 3.42\,nN at 12.1\,mW, 13.65\,nN at 68.4\,mW, 30.13\,nN at
189\,mW, and 40.48\,nN at 278\,mW, confirming the $F/P \propto 1/\sqrt{V}$
throttle law within 2.5\% over 100--300\,V (the 350\,V point sits 3\%
below the ideal, the predicted float tax), with beam escape $\ge 96\%$ and
a floating potential within the 50\,V benign design limit at every point.
The frontier runs discriminate between
pre-registered forms of the ambient-electron collection law: the linear and
square-root laws are refuted and the sub-linear form
$I \propto (1+e\varphi/kT_e)^{\alpha}$ with $\alpha \approx 0.82$--0.89
survives, which fixes the charging tax and places every gated frontier
point within 24\% of a parameter-free ideal-beam power bound -- set by
electron mass and charge alone -- and within 7\% once the two measured
geometry constants are applied; extrapolated one step past its fit range,
the fitted law predicted the 350\,V float to within 1.3\,V (47 predicted,
48.3 measured). A pre-registered run at one-third ambient
density, continued to 2.4\,\textmu s -- the campaign's only fully settled
floating potential -- discriminates the collection law on a second axis:
the square-root form is refuted again, and every fixed-exponent
prediction \emph{overshoots} the measured 42.5\,V float, so the fitted
law is conservative along the density axis. A field-aligned geomagnetic
pair measures a null at flight field strength -- the operating point is
unchanged at 30\,\textmu T -- and a collection tax only at ten times it.
A further pre-registered run varies body
\emph{shape} at fixed drive: tripling the collecting skin while holding the
ram silhouette fixed drops the float from 17.0 to 4.38\,V and \emph{raises}
thrust to 14.22\,nN, because a lower float returns drive energy to the beam.
A final pre-registered pair at 350\,V -- the drive the 400\,km orbit
demands -- closes a $2{\times}2$ voltage--geometry factorial and shows the
two laws \emph{compose}: the slender body floats at 14.0\,V (composed
prediction 11--17\,V) and delivers 43.3\,nN, covering the 400\,km mean
drag with $3.6\times$ charging margin where the compact can sits on the
50\,V limit. This exposes the governing scaling. Drag charges for the ram silhouette while
collection and any body-mounted power supply are paid from the skin, so
thrust demand and
power supply are both areal and vehicle size cancels: the feasibility
condition reduces to the shape ratio (skin\,/\,ram) and the altitude, and is
independent of scale. A minimal executable model calibrated only from the
gated runs, applied to year-long orbit propagations with real solar-activity
and ionosphere data, therefore bounds a corridor that a \O10\,mm body and a
3U CubeSat in end-on flight share identically: at 550--600\,km drag
compensation demands 8--17\,mW mean at 23--29\% duty for the \O10\,mm body;
at 500\,km it remains inside the measured envelope at 41\,mW mean; at
400\,km the mean demand is covered at the 350\,V envelope top, but the
year-round duty (113\%) remains capped by night-side collection at the
compact shape -- the slender geometry is the measured lever, at any size. Supplying the demand is mission design rather than part of the
device -- body-mounted solar cells are the worked example throughout. For
CubeSat-class craft this corresponds to 0.9--3.4\,W and 9--34\,mA, and the
required collection sits only a few times above the bare thermal flux --
markedly less extrapolation than the \O10\,mm anchor itself demands. In an independent predecessor implementation of the same mechanism
(\O14\,mm body, 800\,V) a closed control-volume momentum ledger, validated
on an analytic vacuum fixture and a dedicated full-return null experiment,
established that the recollected return current does not cancel the thrust
(closed thrust $= F_{\mathrm{beam}} - 2.3\%$); in the present campaign the
all-species momentum deposited by collected particles is gated per run and
measures 1.0--8.8\% of beam thrust. We present this as a
feasibility-concept study: the device physics is measured and gated, the
mission projection is model-based with every extrapolated regime explicitly
flagged, and flight validation is future work.
\todo{add lab-experiment sentence once details are in}
\end{abstract}

\section{Introduction}
\label{sec:intro}

Centimetre- to decimetre-class spacecraft in low Earth orbit below
$\sim$600\,km lose altitude to atmospheric drag on timescales of months to a
few years. Extending their operational life requires continuous thrust of a
few to a few tens of nanonewtons, delivered from a power budget of tens to a
few hundred milliwatts -- a regime that sits below where any flight-proven
electric propulsion system is efficient. Cold-gas, PPT, electrospray and FEEP
systems all carry a finite propellant mass in a tank, with feed lines,
valves, or emitter arrays that add mass, cost, and single-point failure modes
disproportionate to a spacecraft this small; gridded-ion and Hall thrusters
are efficient only at powers one to three orders of magnitude above this
budget. The result is a persistent gap: the smallest, cheapest spacecraft
have the least propulsion available to them, at exactly the scale where
station keeping would extend mission life the most cheaply.

We revisit a device long flown in a different role. Every gridded-ion and
Hall thruster carries an electron gun -- the neutralizer -- whose sole job is
to keep the spacecraft from charging as it expels positive ions. We ask what
happens if that same gun, with no ion stage at all, is run as the entire
thruster: fire electrons out through an aperture in an otherwise closed
conducting body, and let the body's own positive charge draw ambient
ionospheric electrons back in to close the circuit. The ionosphere supplies
both the neutralizing return current \emph{and}, in the sense that matters
for a reaction engine, the propellant -- at a mass flow measured in
nanograms per year. No propellant is carried, stored, or metered; the only
consumable is electrical power.

This paper makes five contributions. First, we state the concept's
operating economics in parameter-free form -- an ideal power bound set by
electron mass and charge alone (Sec.~\ref{sec:throttle}) -- and validate
it by measuring the complete thrust--power frontier of the mechanism
across its full hardware voltage range (100--350\,V) with pre-registered,
acceptance-gated particle-in-cell simulation at a realistic ionospheric
operating point (Sec.~\ref{sec:frontier}): the measured points land within
24\% of the bound and deliver ion-thruster-class \emph{energy} conversion
($\eta \approx 0.73$) at a deliberately tiny \emph{momentum} economy
($F/P \approx 0.16$--0.28\,\textmu N/W) -- the pair of numbers that
explains why the device owns the nanonewton regime and nothing above it
(Sec.~\ref{sec:comparison}). Second, the frontier
points discriminate between pre-registered candidate forms of the
ambient-electron collection law (Sec.~\ref{sec:collection}), a
pre-registered one-third-density run -- the campaign's only fully settled
equilibrium -- confirms the discrimination on a second axis and shows the
fitted law is conservative there (Sec.~\ref{sec:density}), and the
350\,V extension lands on the fitted law's out-of-range prediction to
1.3\,V; the surviving
sub-linear law fixes the charging tax, and the measured frontier lands
within 24\% of a parameter-free power bound set by electron mass and charge
alone -- within 7\% once two measured geometry constants are applied
(Sec.~\ref{sec:throttle}). That is the accuracy a feasibility argument
needs, and it is the role the simulations play throughout: proving the
concept against a physics bound, not calibrating a mission-ready thruster.
Third, because the returning electron current is the natural
objection to calling this thrust real, we summarize the closed
control-volume momentum ledger and full-return null experiment of the
predecessor implementation and report the per-run momentum gates of the
present campaign (Sec.~\ref{sec:ledger}). Fourth, a minimal executable
model, calibrated only from the gated runs and with every out-of-envelope
prediction explicitly flagged, converts year-long orbit propagations into a
mission feasibility corridor (Sec.~\ref{sec:application}). Fifth, a further
pre-registered run varies body \emph{shape} rather than drive, and shows the
fitted collection law survives a $3.24\times$ change in collecting area --
and a final pre-registered pair at 350\,V closes a $2{\times}2$
voltage--geometry factorial, showing the voltage and geometry laws
\emph{compose} (slender float 14.0\,V measured against an 11--17\,V
composed prediction); this
licenses a scaling argument in which vehicle size cancels entirely, so the
corridor above applies unchanged to CubeSat-class spacecraft
(Sec.~\ref{sec:slender}) -- the measured body is a \O10\,mm can, but the
result is not about a \O10\,mm can. Throughout we
argue the concept is unusually cheap to test: it needs no new hardware
beyond a high-voltage supply and an existing class of electron gun, and can
retrofit the neutralizer already flying on an ion or Hall thruster
(Sec.~\ref{sec:practical}).

\section{Operating principle}
\label{sec:principle}
The device is a conducting body with a cold-cathode electron gun inside and a
single beam aperture. The full cycle, with measured values from the 200\,V
reference run:
\begin{enumerate}
  \item A high-voltage supply holds the cathode at $-V$ relative to the body;
        the cathode is the most negative node in the system, so emitted
        electrons accelerate \emph{away} from it, through the body aperture.
  \item Electrons leave the body at energy $\approx e(V-\varphi)$ minus
        injection-plane space charge (measured mean exhaust energy
        147.5\,eV at $V = 200$\,V; the measured energy fraction
        $\mathrm{KE}/(V-\varphi) = 0.81$ is constant across the voltage
        range).
  \item Each departed electron leaves the body one charge positive; the body
        potential rises (measured float at the reference: $+16.98$\,V).
  \item The positive body attracts ambient ionospheric electrons, which are
        collected on its conducting surfaces; in equilibrium the collected
        current equals the escaped current exactly (measured current balance
        closes to 1.5--3.5\% at all three operating points). The ionosphere
        is both the propellant reservoir and the neutralizer.
  \item The escaping beam carries net momentum; its reaction is thrust. The
        recollected electrons arrive at thermal-plus-$e\varphi$ energies from
        all directions and deposit a small net momentum (measured and gated
        per run, Sec.~\ref{sec:ledger}).
\end{enumerate}
No valves, no flow control, no gimbals, no plasma production stage: thrust is
throttled by the emission current and supply voltage alone.

\section{The ideal operating law}
\label{sec:throttle}
Before any simulation, the concept's operating economics can be stated with
no free parameters at all. In the ideal limit -- every emitted electron
escapes at its full energy $e(V-\varphi)$, $\varphi \ll V$ -- an electron
beam of current $I$ and per-electron energy $\mathrm{KE}$ carries momentum
flux $F = \sqrt{2 m_e/e}\; I \sqrt{\mathrm{KE}}$; with $F$ in nN, $I$ in mA
and energies in eV the constant is
$\sqrt{2 m_e/e} = 3.37$\,nN/(mA\,$\sqrt{\mathrm{eV}}$) -- electron mass
and charge, nothing else. With ideal beam-supply power $P = VI$,
eliminating the current at a demanded thrust gives the fixed-thrust power
law
\[
P \;=\; \frac{F\,\sqrt{V}}{c_{\mathrm{ideal}}},
\qquad c_{\mathrm{ideal}} = \sqrt{2 m_e/e},
\]
a geometry-free theoretical lower bound. The minimum-power throttle it
implies uses the \emph{largest feasible escaped current} and the
\emph{lowest acceleration voltage} that satisfies the thrust command. The
feasible current is bounded by the emitter, by ambient return-current
availability, by acceptable spacecraft potential, and by successful beam
escape; the concept requires only acknowledging these bounds, not
optimizing against them.

This ideal law is the concept; the rest of the paper is the evidence. The
pre-registered, acceptance-gated simulations that follow ask how close a
real, deliberately non-optimized geometry with real constraints -- beam
space charge, plume divergence, spacecraft charging, finite collection --
comes to the bound (within 20\%,
Sec.~\ref{sec:validation}), whether the law's functional form survives
measurement (the $1/\sqrt{V}$ throttle law, confirmed within 2.5\%,
Sec.~\ref{sec:frontier}), and what sets the one quantity the ideal limit
takes for granted, the return current (Sec.~\ref{sec:collection}).

\section{Simulation method and evidence discipline}
\label{sec:method}
All production runs use WarpX~\cite{warpx22} (v26.5, GPU), RZ electrostatic,
with the spacecraft as an embedded-boundary conductor whose floating
potential is driven each step by the measured charge balance (a
``charge-pump'' update: emitted charge out, collected charge in, body
capacitance from a Gauss-law probe). The anchor deck: a \O10\,mm
$\times$\,5.9\,mm closed can with a 4\,mm exhaust aperture and an internal
cathode disk; ambient plasma at a dayside 400\,km-class row
($n_e = 1.627\times10^{12}$\,m$^{-3}$, $T_e = 1318.8$\,K,
$kT_e = 113.6$\,meV), maintained by a particle reservoir (infinite-ionosphere
recycle -- without refill a floating equilibrium measures reservoir
depletion, not physics); grid $dr = dz = 0.15$\,mm sized to the Debye length
(1.97\,mm); CFL-derived time step; 16 particles per cell per species;
reduced ion mass ($400\,m_e$, caveat in Sec.~\ref{sec:limitations});
800\,ns simulated per run ($\approx$29 capacitive settle times at the
reference point; the density-axis continuation of Sec.~\ref{sec:density}
extends to 2.4\,\textmu s).

Because this paper's claims rest on a small number of expensive runs, the
campaign enforces an explicit evidence contract. Every run freezes and
hashes its full physics config; an atomic manifest records the config hash,
git commit, random seed, code version, and expected-vs-observed final
iteration; interrupted runs are discarded as evidence, never resumed.
Acceptance gates (escape fraction, benign float, steady current balance,
momentum bound, sheath containment, and two independent ledger-vs-particle-dump
charge cross-checks, which close to $\sim$$10^{-9}$) are fixed in versioned
policy files \emph{before} each run; the scaling-law predictions for the
100\,V, 300\,V, and 350\,V points (and the composed prediction for the
slender body at 350\,V), the competing collection-law hypotheses of
Sec.~\ref{sec:collection} (on both the voltage and the density axis),
and the magnetized-axis expectations of Sec.~\ref{sec:magnetized} were
pre-registered in the repository before those runs completed. Below the production stages sits a validation ladder
run under the same contract: an analytic two-electrode Laplace rung, thermal
collection against one-sided flux theory (closes to $\pm$1\%), biased
collection at two potentials, a bare floating body, and electron-gun
emission rungs. The curated metrics, verdicts, frozen configs, and run
manifests of every result used in this paper are committed alongside the
code.

\section{The measured thrust--power frontier}
\label{sec:frontier}

\begin{table}[h]\centering
\begin{tabular}{lccccccc}
\toprule
$V$ [V] & $I$ [mA] & $\varphi$ [V] & $F$ [nN] & $P$ [mW] & $F/P$ [\textmu N/W]
  & escape & KE [eV] \\
\midrule
100 & 0.121 & $+5.40$  & 3.42  & 12.1 & 0.283 & 96.1\% & 77.2 \\
200 & 0.342 & $+16.98$ & 13.65 & 68.4 & 0.200 & 98.4\% & 147.5 \\
300 & 0.630 & $+36.30$ & 30.13 & 189  & 0.159 & 99.0\% & 210.1 \\
350 & 0.793 & $+48.29$ & 40.48 & 278  & 0.146 & 99.1\% & 239.0 \\
\bottomrule
\end{tabular}
\caption{The measured voltage frontier of the \O10\,mm anchor device at the
dayside plasma row, all runs PASS on all pre-registered acceptance gates.
Each point holds the same emission ratio $I/I_{\mathrm{CL}} = 1.46$ (the
space-charge-limited ceiling measured at 200\,V), so beam optics stay
self-similar along the frontier. The 100\,V and 300\,V thrust values landed
on their pre-registered predictions (3.4 and 30\,nN) to better than 1\%,
and the 350\,V extension on its prediction (40.5\,nN) to 0.15\%.
Grid-resolution uncertainty (Sec.~\ref{sec:credibility}): $F \pm 4$\%,
KE $\pm 7$\%, $\varphi \pm 2$\%; particle count, time step, and seed
contribute $\le$0.5\%.}
\label{tab:frontier}
\end{table}

\begin{figure}[h]\centering
\includegraphics[width=\linewidth]{figs/frontier.pdf}
\caption{The measured frontier. Left: thrust along the fixed-perveance path.
Centre: the floating potential against the three pre-registered collection
laws, all anchored at the 200\,V reference; the measured 300\,V float
discriminates them (Sec.~\ref{sec:collection}). The orange square is the
settled-value extrapolation of the 300\,V run's decaying late slope. Right:
thrust per watt against the $1/\sqrt{V}$ law anchored at 200\,V; measured
points agree within 2.5\% over 100--300\,V, with the 350\,V point 3\%
below the ideal -- the growing float tax the two-constant law predicts.}
\label{fig:frontier}
\end{figure}

Two efficiencies describe this device, and the paper keeps them deliberately
side by side. The \emph{energy conversion}
$\eta = P_{\mathrm{jet}}/P_{\mathrm{supply}} \approx 0.73$ at the reference
point -- ion-thruster-class conversion, nearly flat in $V$ -- because the
exhaust is electrons and nearly all of them escape at nearly full energy.
The \emph{impulse economy} $F/P \approx 0.2$\,\textmu N/W at 200\,V,
$\propto 1/\sqrt{V}$ -- roughly 200$\times$ below a gridded ion thruster.
Both are true simultaneously: superb energy-to-jet conversion at tiny
momentum per watt. The second number is why the device only owns the
nanonewton regime; the first is why it deserves to
(Sec.~\ref{sec:comparison}).

The frontier confirms the throttle law $P \propto F\sqrt{V}$ at fixed
thrust: every doubling of voltage costs $\sim$40\% more power, so
\emph{slow-and-many beats fast-and-few} on momentum per watt, and the
cheapest feasible operating point is the lowest voltage whose emission
ceiling still meets the thrust demand. The 100\,V floor run measures
exactly that boundary: escape degrades to 96.1\% (the space-charge tax of
low-voltage optics) but clears the 95\% gate -- the power-optimal end of
the throttle curve is feasible, at 0.283\,\textmu N/W.

\subsection{Validation against the ideal law}
\label{sec:validation}
The frontier measures how close a real, deliberately non-optimized
geometry sits to the bound of Sec.~\ref{sec:throttle}. Writing the
measured beam law as
$F = c_F\, I \sqrt{\kappa(V-\varphi)}$, the fitted thrust slope is
$c_F = 3.27 = 0.97\,c_{\mathrm{ideal}}$ (the 3\% deficit is plume
divergence through the aperture) and the exhaust-energy fraction is
$\kappa = \mathrm{KE}/(V-\varphi) = 0.81$ (injection-plane space-charge
depression, Sec.~\ref{sec:principle}), so the effective constant is
$c_{\mathrm{eff}} = c_F\sqrt{\kappa} = 0.87\,c_{\mathrm{ideal}}$. With
this measured calibration the law reproduces the gated frontier powers to
within 4--7\% at all four operating points (model 11.7 / 65.8 /
177.9 / 258\,mW against measured 12.1 / 68.4 / 189 / 278\,mW), the
residual being
essentially the charging tax $V/(V-\varphi) = 1.06$--$1.16$, which the
measured collection law of Sec.~\ref{sec:collection} supplies when a
tighter number is wanted. The role of the PIC evidence here is to prove
the concept -- a concrete geometry already operates within
$\approx$24\% of the parameter-free physics bound (measured/bound
$= 1.19$--$1.24$ across the frontier, the drift being the growing float
tax), for understood, separable reasons
-- not to certify a mission-ready thruster; the mission
projections of Sec.~\ref{sec:application} use the measured (conservative)
calibration, which is committed with the repository alongside the runs
that produced it.

Off the optimum the bound degrades for known, separable reasons.
Fixed-thrust runs of this geometry at low voltage measured powers
$\approx$1.5--2$\times$ the bound; the excess decomposes into (i) operating
at a non-optimal voltage for the demanded thrust and (ii) genuine
inefficiencies -- beam interception on the enclosed body as low-voltage
optics diverge the plume against the aperture, the sub-unity energy
fraction, and emission-type overheads. The interception loss is specific to
the modeled cathode, gap, and aperture -- the identical 92.4\,V command
transmits 99.99\% in an isolated-gun fixture -- and is not used here as a
universal control law. The attribution is confirmed within the committed
data: the model-selected optimum voltage for the anchor's own 13.65\,nN
demand is 196\,V -- the 200\,V anchor \emph{is} that run -- and there the
bound closes to 4\% with 98.4\% escape, while the
$\approx$1.5--2$\times$ excess appears only at voltages the throttle
principle would never select, where the demand exceeds the emission
ceiling by 2.7--10$\times$. For concept feasibility the ideal bound with
its measured 4--7\% closure is therefore sufficient; closed-loop throttle
optimization for a selected cathode and spacecraft geometry is future
work.

\section{The collection law: a pre-registered discrimination}
\label{sec:collection}
Every escaping electron must be replaced by a collected ambient one; the
floating potential is whatever makes that balance close. The collection
side is therefore the physics that sets the operating point, and its form
was the campaign's principal open question. Writing the collected current
as $I = \beta A\, j_{\mathrm{the}}(n_e,T_e)\,(1+\chi)^{\alpha}$ with
$\chi = e\varphi/kT_e$ and $j_{\mathrm{the}}$ the one-sided thermal flux,
three candidate exponents existed in the project lineage, all agreeing at
the 200\,V anchor ($\chi \approx 150$) and diverging at the 300\,V point
($\chi \approx 320$). Their predictions for the 300\,V float were recorded
in the repository before the run: $\alpha = 1$ (linear, orbital-motion
style) predicts $+31$\,V; $\alpha = 0.82 \pm 0.06$ (fitted across six
equilibria of the predecessor campaign) predicts $+36$\,V; $\alpha = 0.5$
predicts $+90$\,V, which would fail the benign-float gate.

The measured 300\,V float is $+36.30$\,V (tail-averaged under the
pre-registered analysis policy): the linear law is refuted (the float
passes 31\,V and keeps climbing), the square-root law is refuted (nothing
in the trajectory approaches 90\,V), and $\alpha = 0.82$ survives
(Fig.~\ref{fig:frontier}, centre). One caveat was pre-registered with the
test and applies in full: at high $\chi$ the sheath equilibrates on the ion
timescale, and the run's late slope, though decaying ($+44 \to +27$\,mV/ns
over the final 200\,ns), had not reached zero; extrapolating the decay
places the settled float at $\approx$42--48\,V, brushing but not breaking
the 50\,V design limit. A least-squares fit of $\alpha$ across the three
fit points gives $\alpha = 0.89$ (tail-averaged floats) to
$\alpha = 0.85$ (settled-extrapolated floats), bracketing the upper edge of
the pre-registered band.

The fitted law was then tested one step \emph{past} its fit range: a
pre-registered 350\,V run -- the drive the 400\,km orbit demands
($V_{\min} = 304$\,V, Sec.~\ref{sec:application}) -- carried the fitted
law's prediction of $+47$\,V and measured $+48.29$\,V (tail-averaged),
1.3\,V from the prediction. The same settle caveat applies with less
margin: the endpoint float was 51.1\,V and still rising at $+34$\,mV/ns,
so the settled value extrapolates past the 50\,V limit -- at this compact
geometry the 400\,km-enabling drive sits \emph{on} the benign-float limit,
a result the slender geometry of Sec.~\ref{sec:slender} resolves. The
float robs the beam -- the accelerating
potential is always $V - \varphi$ -- so this law is also the charging tax
schedule: $\varphi$ eats 5\% of the supply at 100\,V, 12\% at 300\,V, and
14\% at 350\,V.

\subsection{The density axis: a second discrimination, and a conservative
law}
\label{sec:density}
All three frontier points share one dayside plasma row, so the law's
density dependence -- the $n_e\sqrt{T_e}$ prefactor, and the assumption
that $\alpha$ and $\beta$ do not drift as $r/\lambda_D$ changes -- rested
on theory and on the thermal-collection rung. A pre-registered run at
one-third density ($n_e/3 = 5.42\times10^{11}$\,m$^{-3}$; containment
radius enlarged from 30 to 40.8\,mm for the $\sqrt{3}\times$ larger Debye
length, otherwise the anchor deck unchanged) recorded per-exponent float
predictions before launch: $\alpha = 1$ predicts $+53.4$\,V, $0.893$
predicts $+60.9$\,V, $0.82$ predicts $+68.0$\,V, and $0.5$ predicts
$+160.4$\,V. The 800\,ns run passed every trust gate -- no gross
breakdown of collection in thin plasma -- but its float was still
climbing at run end, so a pre-registered continuation extended the same
run to 2.4\,\textmu s (three times the anchor duration, with the abort
ceiling raised to 180\,V so that no candidate law's prediction could be
clipped by a safety gate).

The continuation delivers the campaign's only fully \emph{settled}
floating potential -- $\varphi = +42.5$\,V with a late slope of
$-0.14$\,V/\textmu s, flat to within noise -- and the verdict surprises
on both ends. The square-root law is refuted a second time,
independently of the voltage axis: the float saturates at 42.5\,V with
the ceiling parked above 160\,V. But every fixed exponent
$\alpha \le 1$ \emph{overshoots}: the enhancement factor $(1+\chi)$ rose
only $2.39\times$ for the $3\times$ density drop, where any fixed
$\alpha \le 1$ requires at least $3\times$. Read as an exponent this is
a two-point secant, $\alpha_{\mathrm{eff}} = \ln 3/\ln 2.39 = 1.26$;
read at the fitted $\alpha$ it is a $\approx$38\% rise in the effective
collection area $\beta A$ -- the direction sheath expansion toward the
orbital-motion limit predicts as $r/\lambda_D$ falls from 2.5 to 1.5.
Operationally the sign is what matters: the fitted law
\emph{over-predicts} the float cost of thin plasma, so the mission
model's density extrapolations (Table~\ref{tab:missions}) err toward
pessimism, and $n_e/3$ does not end the operating envelope at the anchor
drive -- the settled 42.5\,V clears the 50\,V design limit that the
pre-registration expected this run to break. Exhaust energy and thrust
follow the float exactly as the two-constant law demands (measured KE
$= 122.0$\,eV against 122.6 predicted from the injection-plane
potential; $F = 12.39$\,nN; escape 99.1\%). The measurement is one
$3\times$ step, not a sweep: the secant cannot separate exponent drift
from area drift, and densities below $n_e/3$ remain extrapolation --
now with a measured directional bias.

\subsection{Robustness to the geomagnetic field}
\label{sec:magnetized}
Every run above is unmagnetized, while flight is in $\sim$30\,\textmu T.
The only external field orientation compatible with the RZ deck is axial
-- which is precisely the field-aligned-firing flight configuration. A
pre-registered pair added a uniform axial $B_z$ at flight strength
(30\,\textmu T) and at ten times it (300\,\textmu T), with a null
hypothesis bounded in advance at 1$\times$ ($|\Delta\varphi| \le 2$\,V,
$|\Delta F|/F \le 5\%$, $\Delta$escape $\le 1$\,pp) and only the
\emph{direction} of the alternative pre-registered at 10$\times$. At
flight strength the null holds on every bound: float $+17.22$\,V against
the anchor's $+16.98$\,V, thrust 13.64 against 13.65\,nN, escape
unchanged within 0.06\,pp -- the committed operating point is untouched,
as the scale separation predicts (the beam gyroradius, $\ge$0.10\,m, far
exceeds every device dimension, and the Parker--Murphy collection
ceiling sits far above the OML capture radius). At 10$\times$ -- an
amplification instrument, not a flight condition, chosen to push the
thermal-electron gyroradius down to sheath scale -- a collection tax
appears, and it enters entirely through the float: beam formation is
unaffected (escape 98.3\%), but the magnetized skin collects less
effectively, the float rises 32.6\,V over the anchor (still climbing at
run end, so a lower bound on the settled tax), and by
$\mathrm{KE} = \kappa(V-\varphi)$ the exhaust energy falls from 147 to
116\,eV with thrust down 11\%. The two-constant thrust law reproduces
both magnetized points (13.56 and 12.03\,nN predicted against 13.64 and
12.06 measured): $c_F$ is untouched by the field, so the entire effect
is a collection-side tax, consistent with this section's premise that
collection is the physics that sets the operating point. What the pair
does not license is the transverse-field, far-field question -- where
the emitted momentum ends up once the exhaust gyrates at
$r_g \approx 1.4$\,m, four orders of magnitude beyond the simulated
domain -- which cannot be posed in RZ and remains open
(Sec.~\ref{sec:limitations}).

\section{The momentum ledger: is the thrust real?}
\label{sec:ledger}
The first question a plasma physicist asks of this concept: the emitted
current \emph{returns} -- does its momentum cancel the thrust? The answer is
measured at two levels.

In the predecessor implementation (\O14\,mm body, 800\,V, 17.1\,nN at
174\,mW) we closed the full control-volume momentum balance every
simulation window: per-species momentum flux through every boundary,
deposition on the body, emission reaction, field terms
($\int\rho E_z\,dV$, Maxwell stress), and the stored-momentum transient.
Two validation fixtures passed pre-registered gates: a \textbf{vacuum
escape} case with an analytic answer (residual $+0.80$\,nN against a 1\,nN
gate), and a \textbf{full-return null} in which the body is pinned above
the beam energy so escape is 0.0\% -- the impact channel alone reads
$-12.1$\,nN, and the closed thrust is $-0.39\,\text{nN} \approx 0$:
emission recoil, field pull and impact cancel exactly as momentum
conservation demands. At the operating point the ledger closes at
$F = F_{\mathrm{beam}} - 2.3\%$; notably, the return-current impact channel
is cancelled by the field pull on the body, so the na\"ive subtraction
$F_{\mathrm{beam}} - F_{\mathrm{impact}}$ \emph{underestimates} thrust.
\begin{figure}[h]\centering
\includegraphics[width=0.8\linewidth]{figs/ledger_diagram.pdf}
\caption{Momentum-ledger control volume at the predecessor operating point
(\O14\,mm, 800\,V). Solid arrows are directly measured particle momentum
flux; the dashed arrow is the field force on the craft. The impact channel
is almost entirely cancelled by the field pull, not by the escaping beam --
the two together leave the closed thrust within 2.3\% of
$F_{\mathrm{beam}}$ alone.}
\label{fig:ledger}
\end{figure}

In the present campaign the same question is gated on every production run:
the ledger separates $F_{\mathrm{beam}}$ (escaping-beam reaction) from
$F_{\mathrm{net}}$, the $z$-momentum deposited by \emph{all} species
landing on the craft -- collected ambient electrons and ions plus the
returning beam fraction -- and the acceptance policy bounds their ratio.
Measured: $|F_{\mathrm{net}}|/F_{\mathrm{beam}} = 1.0\%$ at 300\,V
(0.30\,nN against 30.13) and 8.8\% at 100\,V, where the slower beam and
lower escape make the return channel relatively largest. The physics
ceiling is worth stating: at current balance the collected current equals
the escaped current, but collected electrons arrive with
$\sim e\varphi$ of directed energy against the beam's
$e\,\mathrm{KE}$, so even if every collected electron arrived from directly
astern the cancellation could not exceed
$\sqrt{\varphi/\mathrm{KE}} \approx 41\%$ (300\,V values); isotropic
arrival -- electrons fall onto the body from all directions -- reduces it
to the measured percent level.

\section{Numerical credibility}
\label{sec:credibility}
\begin{itemize}
  \item \textbf{Code-to-code trajectory validation}: the present (migrated,
        GPU) implementation reproduces the predecessor implementation's
        converged 200\,V anchor point-by-point: float to 0.05\%, escape
        fraction to 0.001\%, $F_{\mathrm{beam}}$ to 0.09\% over 380 logged
        windows.
  \item \textbf{Collection physics vs.\ theory}: the validation ladder's
        thermal-collection rung closes against one-sided flux theory to
        $\pm$1\%; biased rungs measure the finite-size collection
        enhancement at $\chi = 26$ and 88 before any thruster stage runs.
  \item \textbf{Ledger-vs-dump cross-checks}: the per-step scrape ledger is
        checked against independently written particle dumps for both the
        ambient-electron and beam-escape channels on every gated run;
        agreement is at the $10^{-9}$ level (gate: $2\times10^{-2}$).
  \item \textbf{Determinism}: runs are seeded and reproduce their physics
        under rerun; every number in this paper carries its run manifest
        (config hash, seed, code version, git commit).
  \item \textbf{Convergence}: in the predecessor implementation, time step
        ($\times$2 finer), particles-per-cell ($\times$2), and RNG seed each
        change every reported quantity by $\le$0.5\%; grid refinement from
        0.30 to 0.20\,mm cells shifted thrust $+4$--5\% and exhaust energy
        $+10\%$, reported there as the leading open numerical question. For
        the present anchor deck (0.15\,mm cells -- already finer than
        both predecessor grids) we ran a two-axis convergence pair at the
        200\,V anchor. Doubling particles-per-cell (ambient and beam)
        changes every reported quantity by $\le$0.05\% over the full
        800\,ns record, all gates PASS -- particle statistics are a closed
        axis. Grid refinement ($dx\,0.15\!\to\!0.10$\,mm, CFL-locked
        $\Delta t$) shifts $F_{\mathrm{beam}}$ $+4.0$\%, exhaust KE
        $+7.4$\%, and $\varphi$ $-1.8$\%, compared over a matched
        430--536\,ns window (the fine-grid run's step cap ends it at
        536\,ns; the comparison uses the identical physical window on both
        runs, the same fixed-protocol approach as the predecessor's study).
        Grid resolution therefore remains the leading numerical
        uncertainty; we carry it as the error band on
        Table~\ref{tab:frontier} ($F\ {\pm}4$\%, KE ${\pm}7$\%,
        $\varphi\ {\pm}2$\%) and note its sign is conservative for every
        gate: finer grids give \emph{more} thrust and a \emph{lower}
        float. A third grid point to establish the asymptotic order is
        future work.
\end{itemize}

\section{Comparison with existing micropropulsion}
\label{sec:comparison}

\subsection{Related work}
\label{sec:related}
The idea that a spacecraft can dispense with stored propellant by exchanging
charge with its ambient plasma environment has deep roots, though not, to our
knowledge, in exactly the form proposed here. The theoretical foundation for
our device's collection physics is the bare-tether electron-collection theory
of Sanmart\'in, Mart\'inez-S\'anchez, and Ahedo~\cite{san93}, who showed that
an uninsulated conductor floating positive relative to the ambient ionosphere
collects ambient electrons in the orbital-motion-limited (OML) regime -- the
regime whose sub-linear finite-body form our frontier runs measure directly
(Sec.~\ref{sec:collection}); collection by a positively-biased tether under
a \emph{flowing} (ram/drift), rather than stationary, LEO plasma was
subsequently measured experimentally by Siguier et al.~\cite{sig13}, directly
relevant to the ram-wake limitation flagged in Sec.~\ref{sec:limitations}.
This theory was flown and confirmed by NASA's Plasma Motor Generator (PMG)
tether experiment in 1993~\cite{pmg93,kat95}, which demonstrated
hollow-cathode plasma contactors providing a low-impedance
electron-emission-and-collection connection to the ionosphere at each end of
a 500-m tether -- a two-terminal flown precedent for the single-body electron
exchange our device performs internally between its beam aperture and its
own conducting surface. Closer in spirit to a standalone propellantless
thruster is Janhunen's electric solar wind sail
(E-sail)~\cite{jan04}, in which a continuously operating electron gun holds a
tether at high positive potential, actuating tether voltage by ``pump[ing]
out negative charge from the system''~\cite{jan14}, so that it electrostatically
deflects ambient solar-wind protons and converts that deflection into thrust
by momentum conservation; Janhunen separately proposed a LEO-ionospheric
variant, the electrostatic ``plasma brake''~\cite{jan10,jan09neg}, which uses a
tether of the \emph{opposite}, negative polarity to produce Coulomb drag
against the ionospheric ram flow for deorbiting, rather than the positive
polarity of the E-sail or of our device's body. Our device differs from both
in its thrust mechanism -- direct recoil from an escaping electron beam
recollected on a closed body, rather than electrostatic deflection of
ambient ions around a charged tether of either polarity -- but shares their
defining feature that continuous electron emission, not stored propellant,
is what makes the mechanism propellantless. An electrodynamic tether is
also the natural efficiency benchmark: per watt, Lorentz-force thrust on a
kilometre-scale tether far exceeds any reaction engine, ours included. The
comparison inverts for the small spacecraft considered here: neither a
centimetre-class body nor a 3U CubeSat can deploy, stabilize, or survive a
conductor $10^{3}$--$10^{6}$ times its own length,
tether thrust is constrained to the $I\mathbf{L}\times\mathbf{B}$
direction, and a bare tether still requires electron collection and
emission at its terminations -- the same contactor physics this device
\emph{is}, minus the tether. Zubrin's dipole
drive~\cite{zub17} similarly treats ambient space plasma itself as reaction
mass, accelerating collected ions between charged screens while rejecting
electrons, though it targets magnetospheric and interstellar-precursor
regimes rather than low-Earth-orbit drag makeup. A body of ``ambient
atmosphere ion thruster'' patent literature~\cite{dre07} and the broader field
of Atmosphere-Breathing Electric Propulsion (ABEP)~\cite{zhe20} use similar
language -- collecting ambient constituents in place of a propellant tank --
but a materially different mechanism: they collect neutral atmosphere or ram
ions and ionize/accelerate them with a conventional electric-propulsion
stage, retaining an ionization stage, an acceleration stage, and a separate
neutralizer; our device requires none of these, since the ionosphere supplies
its reaction mass pre-ionized and its own neutralizing return current through
the same aperture that emits the beam. A recent general review of
propellantless propulsion concepts~\cite{kez25}, covering gravity assists,
solar/magnetic/electric sails, and quantum-vacuum thrust, does not treat
ambient-ionospheric-electron-exchange thrusters as a category, suggesting the
specific mechanism proposed here -- a sealed body, a cold electron gun, and
nothing else -- has not previously been proposed or evaluated as a standalone
low-Earth-orbit micropropulsion device.

\subsection{Quantitative comparison}
\label{sec:comparison-table}
Table~\ref{tab:comparison} positions our device against eight representative,
flight-proven or flight-qualified micropropulsion systems spanning cold gas
through Hall-effect propulsion, plus a passive photon-pressure baseline, each
sourced from a manufacturer datasheet, a peer-reviewed paper, or NASA's 2026
\emph{State of the Art of Small Spacecraft Technology} report~\cite{nasasoa26}.
The comparison makes three points. First, nanonewton-class \emph{control} is
flight-proven, but nanonewton-class \emph{station-keeping} is not. The
ST7-DRS colloid thrusters on LISA Pathfinder demonstrated 0.1\,\textmu N
(100\,nN) thrust resolution with $\sim$100\,nN/$\sqrt{\mathrm{Hz}}$ noise
against free-floating test masses~\cite{anderson18}, and gallium NanoFEEP
flew sub-micronewton thrust magnitudes on the 1U CubeSat UWE-4 from a few
hundred milliwatts~\cite{bock18} -- so neither a thrust floor nor a power
floor separates this device from the flown state of the art (Enpulsion-class
FEEP states that ``the required power to operate the IFM starts at around
8\,W''~\cite{enpifm}, but that is not the floor of its class). What does
separate it is in the table's hardware column: every incumbent spends a
finite stored propellant through tanks, feed systems, or precision emitters,
so its total impulse is a consumable; the sub-\textmu N precision systems
are drag-free science instruments whose cost and supporting hardware sit
far outside a small-spacecraft budget; and at the gram scale measured here
even NanoFEEP's 160\,g thruster outweighs the vehicle. A verification of
the flight-proven sub-5\,\textmu N landscape against primary sources is
archived with this paper's repository
(\texttt{NANONEWTON\_THRUSTER\_REVIEW.md}). Second, where an explicit or computable jet
efficiency exists, it clusters at 28--45\% for
electrospray~\cite{busekbetmax21}, $\approx$35\% for FEEP~\cite{krejci19},
and $\approx$31--44\% for Hall thrusters~\cite{busekbht200}; our measured
energy conversion (68\% at the predecessor point, 73\% at the anchor
reference) exceeds every value found. Third -- and this is a boundary, not
a boast -- our thrust-per-watt is $\sim$200$\times$ below gridded ion: the
efficiency-indifference boundary is real. The $F/P$ penalty is invisible at
nanonewtons (milliwatts), noticeable at micronewtons ($\sim$10\,W for
CubeSat drag -- disqualifying), and decisive at millinewtons. The natural
handoff -- this device below $\sim$0.1\,\textmu N, electrospray at
micronewtons, ion/Hall at millinewtons -- is part of the claim, not a
concession. ``Continuous'' thrust at
nanonewtons is not a propellant-mass argument (five years of nanonewton
thrust is $\sim$0.3\,g of electrospray propellant) but a \emph{dry-system
floor and total-impulse} argument: tank + feed + valves + PPU is
$\sim$0.3--1\,kg -- a substantial fraction of a kilogram-class CubeSat's
total mass before the first newton-second, and the impulse it buys is
finite -- while this concept's floor is a cathode, a boost converter, and
the spacecraft's own skin, with no total-impulse limit.

\begin{sidewaystable}
\centering\footnotesize
\begin{tabular}{p{2.1cm}p{3.0cm}p{1.9cm}p{1.7cm}p{1.5cm}p{1.3cm}p{3.0cm}p{1.5cm}}
\toprule
System & Example & Thrust & Power & Isp & Energy eff. & Hardware & Ref.\\
\midrule
\textbf{This work} & Ambient-electron exchange & 3.4--43\,nN (\O10\,mm, 100--350\,V); 17\,nN (\O14\,mm, 800\,V) & 12--278\,mW; 174\,mW & $0.7$--$1.5{\times}10^6$\,s & \textbf{68--73\%} ($F/P$ 0.15--0.28\,\textmu N/W) & HV supply + cold gun + aperture. No tank, valves, feed, neutralizer, gimbal & -- \\
Cold gas & GomSpace NanoProp 6U (GOMX-4B)$^\ddagger$ & 0.1--10\,mN & $<$2\,W & 60--110\,s & n/s & Dual tanks, fill valves, filters, 2--3 isolation valves/line, heaters, 4 nozzles & \cite{nasasoa26,gomx4eop} \\
PPT & Mars Space PPTCUP (qualified)$^\ddagger$ & 40\,\textmu N avg & $\sim$2\,W & 655\,s & n/s ($<$20\% typ.) & Cu--W chamber, HV cap bank, EMI box; no tank/valves & \cite{cia15,cia16,nasasoa26} \\
Electrospray & Busek CMNT (LISA Pathfinder) & 4$\times$20\,\textmu N (0.1\,\textmu N resolution) & 16.5\,W & 225\,s & 28/45\%$^\dagger$ (BET-300-P) & Propellant-less CNT neutralizer, low-pressure tanks, PPU & \cite{nasasoa26,busekbetmax21,anderson18} \\
FEEP (1U) & TU Dresden / Morpheus NanoFEEP (UWE-4) & 0--8\,\textmu N cont.\ (peaks 22) & $\sim$0.2--0.3\,W & up to 6000\,s & n/s & Gallium liquid-metal needle heads ($<$6\,g each), propellant heater, CNT neutralizer & \cite{bock18} \\
FEEP & Enpulsion IFM Nano (Flock 3P) & 10\,\textmu N--0.5\,mN & \textbf{8--40\,W} & 2000--5000\,s & $\approx$35\%$^\dagger$ & 28-needle emitter, 2 redundant neutralizer cathodes, PPU & \cite{enpifm,krejci19} \\
RF ion & Busek BIT-3 (Artemis I)$^\ddagger$ & 0.66--1.24\,mN & 56--80\,W & 1400--2640\,s & n/s (PPU $\sim$90\%$^*$) & Grid ion optics, RF cathode neutralizer, iodine feed valve, PPU, 2-axis gimbal & \cite{busekbit3,nasasoa26,nasalunahmap23} \\
Hall & Busek BHT-200 (TacSat-2) & 13\,mN & 200--250\,W & 1390\,s & $\approx$44\%$^\dagger$ & Magnetic discharge channel, hollow cathode, xenon feed, PPU & \cite{busekbht200,nasasoa26} \\
Solar pressure & Passive 10$\times$10\,cm face (computed) & $\approx$91\,nN$^\dagger$ & -- (passive) & -- & -- & None; no throttling, no eclipse operation & \cite{nasatsi} \\
\bottomrule
\end{tabular}
\caption{Comparison with representative flight micropropulsion systems and a
passive baseline. Energy efficiency and thrust-per-watt are stated together
for this work because they diverge by design (see text). $\dagger$ =
derived/computed from other reported figures, not directly quoted.
$\ddagger$ = flight-status caveat: GOMX-4B numbers are flight-demonstrated;
PPTCUP is flight-\emph{qualified} but not confirmed flown as of
\protect\cite{nasasoa26}; BIT-3 flew on Artemis~I but its LunaH-Map flight
demonstration was compromised by a stuck propulsion valve that never freed,
so its in-space performance was not confirmed. $*$ = RF power-supply
electrical efficiency, not overall jet efficiency. n/s = not stated in any
primary source located. Full sourcing and per-number verification notes are
archived alongside this paper's repository.}
\label{tab:comparison}
\end{sidewaystable}

\section{Application: the mission feasibility corridor}
\label{sec:application}
The gated frontier runs calibrate a minimal executable model (committed with
the repository) that applies the throttle principle of
Sec.~\ref{sec:throttle} to every row of a year-long orbit propagation: per
five-minute timestep, the local plasma density and temperature (IRI-2020)
and drag (NRLMSISE-00, real F10.7/Ap, TudatPy
propagation~\cite{tudat}) enter, and the model solves for the minimum-power
operating point $(V, I, \varphi)$ self-consistently with the measured
collection law, the power, and -- critically
-- a per-row envelope flag: any row whose density, $\chi$, or perveance
falls outside what a committed run has measured is counted as
\emph{extrapolated}, never silently averaged in. Five missions
(400--600\,km): Table~\ref{tab:missions}, Fig.~\ref{fig:missions}. The
absolute forces there are quoted for the measured \O10\,mm anchor body; by
the areal argument of Sec.~\ref{sec:scalefree} the \emph{closure} they
express carries to any vehicle of the same shape ratio, CubeSats included.

\begin{table}[h]\centering\footnotesize
\begin{tabular}{lcccccc}
\toprule
altitude / pose & drag mean/max [nN] & feasible & duty cycle & $P$ mean [mW]
 & $\varphi{>}50$\,V & in envelope \\
\midrule
400 km axial   & 32.9 / 92.4 & 36\% & 113\% & 175 & 49\% & 18\% \\
400 km lateral & 21.6 / 60.7 & 63\% & 74\%  & 129 & 33\% & 38\% \\
500 km axial   & 7.6 / 28.4  & 76\% & 39\%  & 41  & 24\% & 29\% \\
550 km axial   & 3.8 / 16.3  & 85\% & 29\%  & 17  & 15\% & 8\%  \\
600 km axial   & 2.0 / 9.6   & 93\% & 23\%  & 8   & 7\%  & 1\%  \\
\bottomrule
\end{tabular}
\caption{Year-long mission sweeps under the throttle principle at the
measured 100--350\,V envelope. ``Feasible'' = a
benign equilibrium exists at the demanded thrust; ``duty cycle'' =
mean drag / mean benign capability ($>$100\% cannot close); ``in
envelope'' = row interpolates committed measurements on every axis.}
\label{tab:missions}
\end{table}

\begin{figure}[h]\centering
\includegraphics[width=0.95\linewidth]{figs/missions.pdf}
\caption{The feasibility corridor. Left: duty cycle needed to cancel mean
drag; above the dashed line the impulse budget cannot close. Right: mean
power at demanded thrust; the body-mounted solar harvest line is an example
supply, not a device requirement.}
\label{fig:missions}
\end{figure}

The corridor reads directly as power demand: at \textbf{550--600\,km} drag
compensation costs 8--17\,mW mean at 23--29\% duty cycles; at
\textbf{500\,km} the impulse budget closes at 39\% duty and the mean demand
is 41\,mW -- the per-row minimum-power solution automatically thrusts
preferentially when the plasma is dense and the charging tax low; at
\textbf{400\,km the mean demand is now inside the measured envelope}
(the 350\,V pair of Secs.~\ref{sec:collection} and \ref{sec:slender}
covers the 32.9\,nN axial mean directly), yet the year-round duty is
113\% for this compact body: the binding constraint has moved from the
voltage ceiling to the benign-float collection cap on thin night-side
rows, where the demanded current drives $\varphi$ past 50\,V -- and we
state that plainly rather than average around it. Supplying these demands is mission design,
not part of the device; for scale, body-mounted cells on this body harvest
roughly 10--30\,mW depending on cell coverage. Two
honesty flags travel with the table. First, the corridor's boundary
behaviour is set by the night-side rows -- thin, cold plasma where the
collection stiffness $\propto 1/(n_e\sqrt{T_e})$ drives the float toward
the 50\,V limit. The density axis now carries one settled measurement
(Sec.~\ref{sec:density}): at one-third density the fitted law
\emph{over-predicts} the float by 10--25\,V, so the corridor's
night-side edge is biased conservative -- the model retires rows the
device could in fact serve. Rows below $n_e/3$, down to the
$\sim$$2\times10^{11}$\,m$^{-3}$ night minimum, remain extrapolation,
now with a measured directional bias rather than none. Second, power
here is beam power $I\,V$; converter and cathode overheads belong to the
flight design and are not modelled.

\subsection{The geometry lever: slender bodies}
\label{sec:slender}
The mission table above holds the body shape fixed, but drag and collection
scale with \emph{different} areas: drag buys the ram silhouette, while
collection and any body-mounted power supply buy the total skin. A slender body flying
end-on decouples them. A \O10\,mm\,$\times$\,100\,mm cylinder presents the
same end cap as the anchor body -- the same nanonewton-class drag -- while its
lateral skin ($\sim$31\,cm$^2$, ten times the anchor's total surface)
collects $\sim$46\,\textmu A of bare thermal current before any float
enhancement and carries skin for an order-of-magnitude larger body-mounted
supply (100--200\,mW of solar cells, as one example), at a mass of tens of
grams. The niche
therefore extends up-mass along the rod (or plate) direction, not the cube:
payload, power, and return current grow with length while the drag bill
does not, and 400\,km -- duty-capped for the squat anchor in
Table~\ref{tab:missions} -- reopens for slender geometries on both the
charging and the power side.

The concern this raised was that the collection exponent is provably
geometry-dependent. OML theory gives $\alpha = 1$ for a sphere and
$\alpha = 0.5$ for a long cylinder; our squat can measured 0.82--0.89,
between the limits, and a slender body should slide toward the cylindrical
value -- weaker enhancement per volt of float, so a hidden charging tax that
the area argument conceals. We therefore pre-registered two competing
hypotheses and measured them: area-only scaling (the can's fitted $\alpha$
holds, so tripling the skin drops the enhancement demand and the float lands
at 4--5\,V) against a cylinder-limit lateral wall (float in the tens of
volts, possibly past the benign gate).

The measurement settles it. A \O10\,mm\,$\times$\,30.5\,mm can
($L/r = 6$, skin $3.24\times$ the anchor's) at the identical 200\,V drive,
commanded current, plasma row, grid and seed floats to
$\phi = 4.38$\,V tail-averaged (4.79\,V at run end, still rising at
$+4.4$\,mV/ns, so $\sim$5--6\,V settled), against the area arithmetic's
4.14--4.66\,V. Hypothesis A is confirmed and the cylinder-limit hypothesis
refuted by an order of magnitude: \textbf{the fitted exponent survives a
$3.24\times$ area change and an aspect-ratio change from $L/r = 0.6$ to 6.}
All six required gates pass; escape holds at 98.4\,\%.

Better still, elongation is not paid for in thrust. Because
$\mathrm{KE} = \kappa(V-\phi)$, a body that floats $3.9\times$ lower keeps
more of its drive: exhaust energy rises from $\sim$147 to 159.7\,eV and beam
thrust from 13.6 to 14.22\,nN, at the same commanded current and the same
ram silhouette. Elongation is free on the thrust side and strictly
favourable on the charging side.

The lever was then tested where it matters most. The 350\,V squat-can run
of Sec.~\ref{sec:collection} left the 400\,km-enabling drive sitting
\emph{on} the 50\,V charging limit, so a final pre-registered run drove
the slender body at the identical 350\,V\,/\,0.793\,mA operating point --
the fourth corner of a $2{\times}2$ voltage--geometry factorial whose
other three corners were already measured. Composing the two fitted laws
predicted a float of 11--17\,V and a thrust of 42--43\,nN; the run
measured $\varphi = 14.00$\,V (tail-averaged; 15.8\,V at run end) and
$F = 43.33$\,nN, with escape at 99.1\%. \textbf{The voltage frontier and
the geometry collection law compose}: together with the anchor they form
a validated two-dimensional operating map over $100$--$350$\,V and
$L/r = 0.6$--$6$, not two independent one-dimensional curves. The mission
reading is direct -- at the drive the 400\,km orbit demands, the shape a
real vehicle takes anyway covers the 32.9\,nN mean drag at $\sim$76\%
duty with a $3.6\times$ margin on the charging limit, and a 7\% thrust
bonus over the compact can from the returned float energy.

Two design trades still travel with the lever: solar cells and bare
collector compete for the same wall area (clad dielectric surfaces do not
collect), and end-on flight is not attitude-free (gravity gradient prefers
the long axis nadir-pointed, which is side-on to ram). A third rule is fixed
by the measured gun physics: the emission ceiling scales as $1/d^{2}$
(Sec.~\ref{sec:frontier}), so a slender body must keep the
cathode--aperture gap at its short design value and grow the body
\emph{around} the gun -- an emitter placed at the far end of a long can
starves and scatters its own beam before the slit. We enforced this with a
cathode standoff that pins the gap while the can grows; the run's
demand-to-ceiling ratio, 1.457, matches the anchor's 1.46.

What remains extrapolation is aspect ratio \emph{beyond} $L/r = 6$, where
the cylinder limit must eventually bite, and the $r \lesssim \lambda_D$ end
that converges to the bare-tether collection law of
Sanmart\'in et al.~\cite{san93} grounding our related work.

\subsection{Size cancels: the feasibility condition is scale-free}
\label{sec:scalefree}
The geometry measurement licenses a scaling argument that changes what this
device is for. Drag charges for the ram silhouette $A_{\mathrm{ram}}$, so the
thrust demand, the beam current and hence the input power are all
proportional to $A_{\mathrm{ram}}$; power \emph{per unit ram area} is
therefore a property of altitude and drive voltage alone
(Table~\ref{tab:scalefree}, top). Any body-mounted power supply is likewise
areal on the skin -- solar cells, the worked example here, deliver a fixed
power per unit cell area whatever the vehicle. \textbf{Vehicle size
cancels from both sides of the power balance}, and what survives is the shape
ratio $A_{\mathrm{skin}}/A_{\mathrm{ram}}$ together with the altitude.

\begin{table}[t]\centering\small
\caption{The power balance is areal on both sides, so feasibility is set by
shape and altitude, not by size. Margin = example supply\,/\,demand at
100\,V; the example supply is body-mounted solar, 30\% cells over one third
of the skin at 25\% orbit-average illumination (102\,W/m$^2$ of cell area) --
the size-free ratio is the result, not the absolute number. Generated by
\texttt{model/scale\_analysis.py}.}
\label{tab:scalefree}
\begin{tabular}{lrrrr}
\toprule
altitude & 400\,km & 500\,km & 550\,km & 600\,km \\
\midrule
power demand per m$^2$ of ram & 1451\,W & 335\,W & 170\,W & 88\,W \\
\midrule
\multicolumn{5}{l}{\emph{margin, by body (skin/ram ratio in parentheses)}}\\
\O10\,mm can, squat \hfill (4) & 0.1 & 0.4 & 0.8 & 1.6 \\
1U cube, face-on \hfill (6) & 0.1 & 0.6 & 1.2 & 2.3 \\
\O10\,mm can, slender \hfill (14) & 0.3 & \textbf{1.4} & \textbf{2.8} & \textbf{5.4} \\
3U CubeSat, end-on \hfill (14) & 0.3 & \textbf{1.4} & \textbf{2.8} & \textbf{5.4} \\
6U CubeSat, end-on \hfill (12) & 0.3 & 1.2 & 2.4 & 4.6 \\
\bottomrule
\end{tabular}
\end{table}

The slender \O10\,mm can and the 3U CubeSat return \emph{identical} margins,
because they have the same shape ratio; the 400\,km wall is equally
scale-free and closes for nothing. Read the other way, the corridor of
Table~\ref{tab:missions} is not a small-body result at all -- it is a
statement about aspect ratio and altitude that a kilogram-class,
payload-carrying CubeSat inherits unchanged. In absolute terms a 3U in end-on flight needs 8.8\,mA at 0.88\,W
to hold 600\,km, 17\,mA at 1.7\,W at 550\,km, and 33\,mA at 3.4\,W at
500\,km: values a CubeSat power system can supply.

Two things become \emph{easier} with size. First, larger skins collect more,
so the enhancement demanded over bare thermal collection falls -- a 3U needs
$4.3\times$ the thermal flux at 600\,km where the \O10\,mm anchor needs
$15\times$, and the anchor frontier runs at $\chi \approx 150$--320. Since
bare thermal collection is the rung validated to $\pm1\%$ and the enhancement
exponent is the fitted quantity, \textbf{larger bodies rest on less
extrapolation, not more}. Second, lower $\chi$ means a lower float, which by
$\mathrm{KE} = \kappa(V-\varphi)$ returns drive energy to the beam.

The honest counterweight is a regime change we have not measured. Every
committed run sits at $r/\lambda_D \approx 2.5$; CubeSat radii are
25--60\,$\lambda_D$, where the orbital-motion picture does not apply at all
and collection proceeds through a thin sheath that grows with $\varphi$.
Estimating with that (Child-sheath) model instead of the fitted exponent puts
the 3U float at $\sim$12\,V at 600\,km, $\sim$25\,V at 550\,km and
$\sim$47\,V at 500\,km -- benign in the corridor, tightening at its lower
edge, and past the benign gate for the largest bodies at 500\,km. That is an
estimate, not a calibration. A single large-body run would settle it, and it
is cheaper than the runs already reported, because a lower-$\chi$ equilibrium
settles faster.

\section{Practical realization}
\label{sec:practical}
\subsection{Bill of materials}
A conducting structure, a milliwatt-class high-voltage supply (100--350\,V
for the measured design; up to 800\,V demonstrated in simulation at larger
scale), a cold electron gun with its cathode at $-V$, and an aperture. No
tank, valves, feed lines, RF stages, or gimbals. The dry-system floor is a
cathode, a boost converter, and the spacecraft's own skin -- grams of added
hardware on a kilogram-class CubeSat, where every tanked system starts at
0.3--1\,kg of dry mass before the first newton-second.
\todo{cost estimate}
\subsection{Retrofit path}
Gridded-ion and Hall thrusters already fly hollow-cathode/filament electron
sources as neutralizers. After propellant exhaustion, the neutralizer plus
the spacecraft structure \emph{is} this thruster: a software-and-wiring
change buys drag-makeup-class thrust for mission life extension.
\subsection{Laboratory demonstration}
\todo{user's vacuum-chamber experiment: negative-cathode electron
acceleration demonstrated on bench hardware -- geometry, voltages, pressure,
measured currents, photos. Details pending.}

\section{Limitations and future work}
\label{sec:limitations}
We state plainly where this study is known to be incomplete, so a reader can
judge exactly how far the present results extend. This is a
feasibility-concept paper: the device physics is measured under the evidence
contract of Sec.~\ref{sec:method}; the mission projection is model-based
with flagged extrapolation; full mission validation is future work.

\begin{itemize}
  \item \textbf{The collection law's density axis is measured at one
        $3\times$ step.} A pre-registered one-third-density run, continued
        to a settled equilibrium, shows the fitted law is
        \emph{conservative} there -- it over-predicts the float
        (Sec.~\ref{sec:density}). The step is a two-point secant: exponent
        drift ($\alpha_{\mathrm{eff}} = 1.26$) cannot be separated from
        area drift ($\beta A$ up $\approx$38\%) with one step, and
        densities below $n_e/3$ -- the deepest night rows -- remain
        extrapolation, with a measured directional bias. The model still
        flags every such row (Table~\ref{tab:missions}).
  \item \textbf{The collection law's geometry axis is measured at two
        shapes, not mapped.} The exponent is geometry-dependent by theory
        (sphere 1, long cylinder 0.5; measured 0.82--0.89 for the can), and
        the pre-registered slender runs (Sec.~\ref{sec:slender}) confirm it
        survives to $L/r = 6$ and $3.24\times$ skin area at both 200 and
        350\,V -- a $2{\times}2$ factorial, not a continuous map. Aspect
        ratios beyond that remain extrapolation resting on the theory-safe
        thermal-collection floor plus flags, not on the fitted law.
  \item \textbf{No run sits in the CubeSat sheath regime.} All committed
        bodies are at $r/\lambda_D \approx 2.5$, where orbital-motion
        collection applies. CubeSat radii are 25--60\,$\lambda_D$: a thin,
        $\varphi$-growing sheath instead, for which we quote a Child-sheath
        \emph{estimate} (Sec.~\ref{sec:scalefree}) and not a calibration.
        The scale-free power balance of Table~\ref{tab:scalefree} does not
        depend on this -- it is areal bookkeeping -- but the floating
        potentials of CubeSat-class bodies do. This is the single highest-value
        follow-on measurement, and the cheapest, since lower-$\chi$
        equilibria settle faster than the runs already reported.
  \item \textbf{Grid resolution is the leading numerical uncertainty.}
        Particle count is converged to $\le$0.05\% (this work) and time
        step and seed to $\le$0.5\% (predecessor), but the measured
        0.15$\to$0.10\,mm grid shift is $+4.0$\% on thrust and $+7.4$\% on
        exhaust energy (matched-window protocol,
        Sec.~\ref{sec:credibility}), carried as the error band on every
        frontier number. Two grid points cannot establish an asymptotic
        order; a third is future work. The shift's sign is conservative
        for all acceptance gates.
  \item \textbf{Finite-time equilibrium on the ion clock.} The 300\,V float
        was still relaxing at run end (Sec.~\ref{sec:collection}); the
        settled-value band (42--48\,V) is carried through the $\alpha$ fit
        rather than hidden, and brushes the 50\,V design limit. The
        350\,V squat float is the sharpest case: gated at 48.3\,V on the
        tail mean but 51.1\,V and still rising at run end, so its settled
        value likely exceeds the limit -- the compact body sits on the
        boundary, and the slender body (14--17\,V band) carries the
        margin. The
        2.4\,\textmu s density-axis continuation is the only fully settled
        float (Sec.~\ref{sec:density}) and measures the settle time for
        its row at $\sim$2\,\textmu s; the anchor-row floats remain
        800\,ns reads with disclosed slopes.
  \item \textbf{Reduced ion mass.} Simulations use $400\,m_e$ rather than
        real O$^+$ ($29{,}166\,m_e$) to keep ion timescales affordable; the
        sensitivity of the float to ion mass has not been measured.
  \item \textbf{Stationary ambient plasma.} No bulk flow is modelled. For
        electron collection the ram anisotropy is small
        ($v_{\mathrm{orb}}/v_{\mathrm{th},e} \approx 5\%$); for a body
        flying end-on the ram-ion wake is the largest unmodelled effect and
        is deferred, not estimated.
  \item \textbf{Two-dimensional axisymmetry; transverse geomagnetic field
        unmodelled.} The field-aligned (axial) component is measured: a
        null at flight strength and a collection tax only at $10\times$
        (Sec.~\ref{sec:magnetized}). The transverse component -- and with
        it the far-field fate of the emitted momentum, which gyrates at
        $r_g \approx 1.4$\,m against a 30\,mm domain -- cannot be posed
        in RZ and is the largest unexamined question.
  \item \textbf{System power.} All powers are beam power $I\,V$. Cathode
        heating (zero for field emitters, nonzero otherwise) and converter
        efficiency are flight-design quantities outside this paper's scope.
  \item \textbf{Throttle optimization is deferred.} Quoted powers are the
        ideal beam-supply bound of Sec.~\ref{sec:throttle}, validated to
        4--7\% in the high-escape regime. Off-optimum low-voltage
        operation in this geometry measured up to
        $\approx$1.5--2$\times$ the bound, an excess attributable to
        non-optimal voltage plus geometry-specific losses (beam
        interception, plume divergence, emission-type overheads). The
        200\,V anchor already sits at the model-selected optimum for its
        own demand and closes to 4\%, confirming the attribution at one
        demand; targeted runs at other demands are cheap follow-ons, and
        closed-loop controller design for a selected cathode and geometry
        is future engineering work.
  \item \textbf{Full mission simulation.} Battery state of charge, seasonal
        density variation, attitude coupling, and a flight-software
        implementation of the throttle are future work; the model's per-row
        accounting (Sec.~\ref{sec:application}) is the input such a study
        needs, not a substitute for it.
\end{itemize}

\section{Conclusion}
\label{sec:conclusion}
Drag compensation for small spacecraft at 450--600\,km requires
nanonewtons to micronewtons of continuous thrust from milliwatts to a few
watts of power. Flight-proven micropropulsion reaches that scale in control
resolution but not in kind: every flown system spends a finite stored
propellant through tanks, feed systems, and precision emitters whose mass,
cost, and total-impulse limit are disproportionate to a spacecraft this
small. We have measured, under a pre-registered evidence
contract, the complete thrust--power frontier of a device that fills the
slot with zero propellant and zero net mass flux: a sealed conducting body,
a cold electron gun, and the ionosphere as both propellant reservoir and
return electrode. Three results support treating this as a credible
propulsion mechanism rather than a numerical curiosity. First, the measured
frontier -- 3.4\,nN at 12\,mW to 40.5\,nN at 278\,mW across the full
hardware voltage range -- confirms the $F/P \propto 1/\sqrt{V}$ throttle
law and delivers ion-thruster-class energy conversion
($\approx$0.73) at a thrust-per-watt $\sim$200$\times$ below gridded ion:
the pair of numbers that assigns this device the nanonewton regime and
nothing above it. Second, the returning current does not cancel the
thrust: a closed momentum ledger with a dedicated full-return null
experiment establishes it in the predecessor implementation, and per-run
momentum gates bound the return channel at 1--9\% of beam thrust in the
present one. Third, the measured collection law -- discriminated on both
the voltage and the density axis, settled on the latter, conservative
where it extrapolates, undisturbed by a flight-strength field-aligned
geomagnetic field, and shown to \emph{compose} with the geometry law
across a fully measured $2{\times}2$ voltage--geometry factorial (the
slender body at the 400\,km-enabling 350\,V drive landed dead-centre of
its composed prediction) -- fixes the charging tax, and
the resulting two-constant analytical power model -- validated to within
7\% of every gated frontier point -- applied to year-long orbit
propagations with every extrapolated regime flagged, bounds a mission
corridor stated as power demand: 8--17\,mW at 550--600\,km, 41\,mW at
500\,km, and at 400\,km a mean demand now covered inside the measured
envelope (with the charging margin belonging to the slender shape) while
the year-round duty stays honestly capped by night-side collection -- a
corridor that, by the measured size-cancellation argument,
a 3U CubeSat in end-on flight inherits unchanged at 0.9--3.4\,W and
9--34\,mA. The device is inexpensive and low-risk to test -- no hardware
beyond a high-voltage supply and an electron gun of a kind that already
flies as a neutralizer -- which suggests both a rapid path to an on-orbit
demonstration and a retrofit use for existing propulsion systems after
their propellant is exhausted.

\begin{thebibliography}{99}

\bibitem{warpx22} J.-L.~Vay et al., ``Modeling of a chain of three plasma
  accelerator stages with the WarpX electromagnetic PIC code on GPUs,''
  \emph{Physics of Plasmas}, vol.~28, 2021, 023105; L.~Fedeli, A.~Huebl,
  et al., ``Pushing the Frontier in the Design of Laser-Based Electron
  Accelerators with Groundbreaking Mesh-Refined Particle-In-Cell
  Simulations on Exascale-Class Supercomputers,'' SC22 (Gordon Bell
  Prize), 2022.

\bibitem{tudat} Tudat Space (TudatPy): TU Delft Astrodynamics Toolbox,
  \url{https://docs.tudat.space}; propagation with NRLMSISE-00 atmosphere
  and IRI-2020 ionosphere.

\bibitem{san93} J.~R. Sanmart\'in, M.~Mart\'inez-S\'anchez, and E.~Ahedo,
  ``Bare Wire Anodes for Electrodynamic Tethers,'' \emph{Journal of
  Propulsion and Power}, vol.~9, no.~3, 1993, pp.~353--360.

\bibitem{pmg93} NASA, ``Plasma Motor Generator (PMG) Electrodynamic Tether
  Experiment,'' flight June 26, 1993, NASA Technical Reports Server, accession
  19960002041.

\bibitem{kat95} I.~Katz et al., ``Plasma turbulence enhanced current
  collection: Results from the plasma motor generator electrodynamic tether
  flight,'' \emph{Journal of Geophysical Research: Space Physics}, 1995,
  \href{https://doi.org/10.1029/94JA03142}{doi:10.1029/94JA03142}.

\bibitem{jan04} P.~Janhunen, ``Electric sail for spacecraft propulsion,''
  \emph{Journal of Propulsion and Power}, vol.~20, no.~4, 2004, pp.~763--764.

\bibitem{jan14} P.~Janhunen, ``Coulomb drag devices: electric solar wind sail
  propulsion and ionospheric deorbiting,'' arXiv:1404.7430 [astro-ph.IM],
  2014.

\bibitem{jan10} P.~Janhunen, ``Electrostatic plasma brake for deorbiting a
  satellite,'' \emph{Journal of Propulsion and Power}, vol.~26, no.~2, 2010,
  pp.~370--372.

\bibitem{jan09neg} P.~Janhunen, ``On the feasibility of a negative polarity
  electric sail,'' \emph{Annales Geophysicae}, vol.~27, 2009, pp.~1439--1447.

\bibitem{sig13} J.-M.~Siguier, P.~Sarrailh, J.-F.~Roussel, V.~Inguimbert,
  G.~Murat, and J.~Sanmart\'in, ``Drifting plasma collection by a positive
  biased tether wire in LEO-like plasma conditions: current measurement and
  plasma diagnostic,'' \emph{IEEE Transactions on Plasma Science}, vol.~41,
  no.~12, 2013, pp.~3380--3386.

\bibitem{zub17} R.~Zubrin, ``Dipole Drive for Space Propulsion,'' \emph{Journal
  of the British Interplanetary Society}, vol.~70, 2017, pp.~442--448.

\bibitem{dre07} G.~A.~Dressler, ``System and method for an ambient atmosphere
  ion thruster,'' US Patent 7,306,189 B2, Northrop Grumman Systems
  Corporation, filed Dec.~11, 2006, granted Dec.~11, 2007.

\bibitem{zhe20} P.~Zheng et al., ``A Comprehensive Review of
  Atmosphere-Breathing Electric Propulsion Systems,'' \emph{International
  Journal of Aerospace Engineering}, 2020, Article ID 8811847.

\bibitem{esaramep} European Space Agency, ``World-first firing of
  air-breathing electric thruster,'' ESA news, March 2018.

\bibitem{kez25} R.~Ya.~Kezerashvili, ``Propellantless space exploration,''
  arXiv:2510.21743, 2025.

\bibitem{nasasoa26} NASA, ``State of the Art of Small Spacecraft Technology,''
  2026 edition, In-Space Propulsion chapter (Ch.~4).

\bibitem{gomx4eop} ESA / eoPortal, ``GomX-4 (GomSpace Express-4)'' satellite
  mission page, \url{https://www.eoportal.org/satellite-missions/gomx-4}.

\bibitem{cia15} S.~Ciaralli, M.~Coletti, and S.~B.~Gabriel, ``Qualification of
  a Pulsed Plasma Thruster for Cubesat Propulsion (PPTCUP),''
  IEPC-2015-208/ISTS-2015-b-208, 34th International Electric Propulsion
  Conference, Kobe-Hyogo, Japan, 2015.

\bibitem{cia16} S.~Ciaralli, M.~Coletti, and S.~B.~Gabriel, ``Results of the
  qualification test campaign of a Pulsed Plasma Thruster for Cubesat
  Propulsion (PPTCUP),'' \emph{Acta Astronautica}, vol.~121, 2016,
  pp.~314--322.

\bibitem{busekelec17} Busek Co.\ Inc., ``All-Electric CubeSat Propulsion
  Technologies,'' 6th International Interplanetary CubeSat Workshop, May~31,
  2017.

\bibitem{busekbetmax21} Busek Co.\ Inc., ``BET-MAX Electrospray Thruster
  System'' datasheet, v1.0, August 2021.

\bibitem{enpifm} Enpulsion GmbH, ``Introduction to the IFM Nano Thruster''
  product overview datasheet.

\bibitem{krejci19} D.~Krejci, B.~Seifert, C.~Scharlemann, et al., ``Full
  Performance Mapping of the IFM Nano Thruster, Including Direct Thrust
  Measurements,'' \emph{Journal of Space Safety Engineering}, vol.~8, no.~2,
  2019, p.~882.

\bibitem{anderson18} G.~Anderson et al., ``Experimental results from the
  ST7 mission on LISA Pathfinder,'' \emph{Physical Review D}, vol.~98,
  102005, 2018.

\bibitem{bock18} D.~Bock and M.~Tajmar, ``Highly miniaturized FEEP
  propulsion system (NanoFEEP) for attitude and orbit control of CubeSats,''
  \emph{Acta Astronautica}, vol.~144, 2018, pp.~422--428; see also
  ``UWE-4: First Electric Propulsion on a 1U CubeSat,'' 2020 (in-orbit
  operation at $\sim$0.2--0.3\,W).

\bibitem{busekbit3} Busek Co.\ Inc., ``BIT-3 RF Ion Thruster'' product page,
  \url{https://www.busek.com/bit3}.

\bibitem{nasalunahmap23} NASA Science, ``NASA's LunaH-Map Mission Ends,
  Validates Science Instrument Performance,'' Aug.~3, 2023.

\bibitem{busekbht200} Busek Co.\ Inc., ``BHT-200'' product page,
  \url{https://www.busek.com/bht200}.

\bibitem{nasatsi} NASA Goddard Space Flight Center, ``Solar Irradiance
  Science,''
  \url{https://earth.gsfc.nasa.gov/climate/projects/solar-irradiance/science}.

\end{thebibliography}

\end{document}
